\Article{16}
Zunächst ergibt sich $\Sigma Tex=T\cos u\, .\, \Sigma ae-T\sin u\, .\, \Sigma be=mT\cos u$, wenn wir das \SourceLeaf{107}Moment des freien Magnetismus der ersten Nadel $\Sigma ae$ mit $m$ bezeichnen, denn es gilt bekanntlich $\Sigma be=0$. Der Beitrag zu $\frac{\mathrm{d}\Omega}{\mathrm{d}u}$ aus dem ersten Term von $\Omega$ ist dann $=-mT\sin u$.

Setzen wir der Kürze halber
\[
\begin{aligned}
k={}&\alpha\cos\psi+\beta\sin\psi+A\cos(\psi-U)+B\sin(\psi-U)\\
&-a\cos(\psi-u)-b\sin(\psi-u)\\
l={}&(\alpha\sin\psi-\beta\cos\psi+A\sin(\psi-U)-B\cos(\psi-U)\\
&-a\sin(\psi-u)+b\cos(\psi-u))^2+(C-c)^2
\end{aligned}
\]
so gilt $rr=(R+k)^2+l$.

Da bei brauchbaren Versuchen $R$ viel größer sein muss als die Abmessungen jeder der beiden Nadeln, lässt sich die Größe $\frac1{r^{n-1}}$ in die rasch konvergierende Reihe
\[
\begin{aligned}
&R^{-(n-1)}-(n-1)kR^{-n}+\left(\frac{nn-n}{2}kk-\frac{n-1}{2}l\right)R^{-(n+1)}\\
&\qquad-\left(\tfrac16(n^3-n)k^3-\tfrac12(nn-1)kl\right)R^{-(n+2)}+\text{usw.}
\end{aligned}
\]
entwickeln. Das Bildungsgesetz der Reihe ließe sich leicht angeben, wenn es der Mühe wert wäre. Die einzelnen Terme der Summe $\Sigma\frac{eE}{r^{n-1}}$, die sich nach Einsetzen der Werte von $k,l$ ergeben, enthalten jeweils einen Faktor der Form
\[
\Sigma eEa^\lambda b^\mu c^\nu A^{\lambda'}B^{\mu'}C^{\nu'}
\]
Dieser Faktor ist gleich dem Produkt der Faktoren $\Sigma ea^\lambda b^\mu c^\nu$, $\Sigma EA^{\lambda'}B^{\mu'}C^{\nu'}$, die vom magnetischen Zustand der ersten beziehungsweise der zweiten Nadel abhängen. Was sich in dieser Hinsicht allgemein feststellen lässt, beschränkt sich auf die Gleichungen
\[
\begin{gathered}
\Sigma e=0,\ \Sigma ea=m,\ \Sigma eb=0,\ \Sigma ec=0,\\
\Sigma E=0,\ \Sigma EA=M,\ \Sigma EB=0,\ \Sigma EC=0
\end{gathered}
\]
wobei wir mit $M$ das Moment des freien Magnetismus der zweiten Nadel bezeichnen. Betrachten wir den besonderen Fall, in dem die Form der ersten Nadel und die Verteilung ihres Magnetismus in Längsrichtung symmetrisch sind. Dabei entsprechen sich stets zwei Elemente, für die $a$ und $e$ entgegengesetzte und $b$ und $c$ gleiche Werte haben, und die Mitte fällt mit dem Punkt $h$ zusammen. Dann gilt stets $\Sigma ea^\lambda b^\mu c^\nu=0$ für jeden geraden Wert der Zahl $\lambda+\mu+\nu$. Entsprechendes gilt für die zweite Nadel, wenn ihre Form und die Verteilung ihres Magnetismus bezüglich des Punkts $H$ symmetrisch sind. Allgemein verschwinden daher in der Summe $\Sigma\frac{eE}{r^{(n-1)}}$ die Koeffizienten der Potenzen $R^{-(n-1)}$ und $R^{-n}$. Im besonderen Fall, in dem beide Nadeln symmetrisch geformt und symmetrisch magnetisiert sind, zugleich die Mitte der ersten Nadel mit $h$ und $h'$ zusammenfällt und die Mitte der zweiten mit $H$, verschwinden auch die Koeffizienten der Potenzen $R^{-(n+2)},R^{-(n+4)},R^{-(n+6)}$ usw. Sind jene Bedingungen annähernd \SourceLeaf{108}erfüllt, so müssen diese Koeffizienten wenigstens sehr klein werden. Der Hauptterm, der sich aus der Entwicklung des zweiten Teils von $\Omega$, nämlich von $-\Sigma\frac{eE}{(n-1)r^{(n-1)}}$, ergibt, ist
\[
\begin{aligned}
&=-\tfrac12 R^{-(n+1)}(n\Sigma eEkk-\Sigma eEl)\\
&=mMR^{-(n+1)}(n\cos(\psi-U)\cos(\psi-u)-\sin(\psi-U)\sin(\psi-u))
\end{aligned}
\]

Daraus folgt, dass sich der Teil von $\frac{\mathrm{d}\Omega}{\mathrm{d}u}$, der der Wirkung der zweiten Nadel entspricht, durch eine Reihe der Form
\[
fR^{-(n+1)}+f'R^{-(n+2)}+f''R^{-(n+3)}+\text{usw.}
\]
ausdrücken lässt. Die Koeffizienten sind rationale Funktionen der Kosinus und Sinus der Winkel $\psi,u,U$ und der Größen $\alpha,\beta$. Sie enthalten außerdem konstante Größen, die vom magnetischen Zustand der Nadeln abhängen. Insbesondere gilt
\[
f=mM(n\cos(\psi-U)\sin(\psi-u)+\sin(\psi-U)\cos(\psi-u))
\]
Die vollständige Entwicklung der folgenden Koeffizienten $f',f''$ usw. ist für unseren Zweck nicht nötig. Folgende Beobachtungen genügen.

1) Im eben beschriebenen Fall vollkommener Symmetrie verschwinden die Koeffizienten $f',f'''$ usw.

2) Vergrößert man bei unveränderten übrigen Größen $\psi$ um zwei rechte Winkel (oder, was dasselbe ist, nimmt man den Abstand $R$ auf derselben rückwärts verlängerten Geraden auf der anderen Seite des Punkts $h'$), so behalten die Koeffizienten $f,f'',f''''$ usw. ihre Werte, während $f',f''',f^{\mathrm{v}}$ usw. die entgegengesetzten Werte annehmen. Die Reihe geht also über in
\[
fR^{-(n+1)}-f'R^{-(n+2)}+f''R^{-(n+3)}-\text{usw.}
\]
Das ergibt sich leicht daraus, dass bei dieser Änderung von $\psi$ die Größe $k$ in $-k$ übergeht und $l$ unverändert bleibt.

